% Custom command definitions for Social Deduction Game role cards
% Adapted from D&D item card commands to match the original aesthetic

% Role card creation command using D&D styling
% \createrolecard{name}{type}{image}{objective}{abilities}{victory}{strategy}{warnings}
\newcommand{\createrolecard}[8]{
    \begin{rolecardauto}{#1}
        \begin{center}
            \begin{tcolorbox}[width=\linewidth, boxrule=3.0pt, colframe=black, boxsep=0pt, top=0pt, bottom=0pt, left=0pt, right=0pt, enhanced, borderline={0.5mm}{0mm}{black}, interior style={fill overzoom image=img/paper.jpg, fill image opacity=1}]
                \includegraphics[width=\linewidth]{#3}
            \end{tcolorbox}
        \end{center}
        
        \begin{lowerpartauto}
            \RoleType{#2}
            \RoleObjective{#4}
            \RoleAbilities{#5}
            \VictoryCondition{#6}
            \StrategyTip{#7}
            \RoleWarnings{#8}
        \end{lowerpartauto}
    \end{rolecardauto}
}

% Town role card with blue coloring
\newcommand{\createtownrolecard}[8]{
    \begin{townroleauto}{#1}
        \begin{center}
            \begin{tcolorbox}[width=\linewidth, boxrule=3.0pt, colframe=black, boxsep=0pt, top=0pt, bottom=0pt, left=0pt, right=0pt, enhanced, borderline={0.5mm}{0mm}{black}, interior style={fill overzoom image=img/paper.jpg, fill image opacity=1}]
                \includegraphics[width=\linewidth]{#3}
            \end{tcolorbox}
        \end{center}
        
        \begin{lowerpartauto}
            \RoleType{#2}
            \RoleObjective{#4}
            \RoleAbilities{#5}
            \VictoryCondition{#6}
            \StrategyTip{#7}
            \RoleWarnings{#8}
        \end{lowerpartauto}
    \end{townroleauto}
}

% Mafia role card with red coloring
\newcommand{\createmafiarolecard}[8]{
    \begin{mafiaroleauto}{#1}
        \begin{center}
            \begin{tcolorbox}[width=\linewidth, boxrule=3.0pt, colframe=black, boxsep=0pt, top=0pt, bottom=0pt, left=0pt, right=0pt, enhanced, borderline={0.5mm}{0mm}{black}, interior style={fill overzoom image=img/paper.jpg, fill image opacity=1}]
                \includegraphics[width=\linewidth]{#3}
            \end{tcolorbox}
        \end{center}
        
        \begin{lowerpartauto}
            \RoleType{#2}
            \RoleObjective{#4}
            \RoleAbilities{#5}
            \VictoryCondition{#6}
            \StrategyTip{#7}
            \RoleWarnings{#8}
        \end{lowerpartauto}
    \end{mafiaroleauto}
}

% Neutral role card with orange coloring
\newcommand{\createneutralrolecard}[8]{
    \begin{neutralroleauto}{#1}
        \begin{center}
            \begin{tcolorbox}[width=\linewidth, boxrule=3.0pt, colframe=black, boxsep=0pt, top=0pt, bottom=0pt, left=0pt, right=0pt, enhanced, borderline={0.5mm}{0mm}{black}, interior style={fill overzoom image=img/paper.jpg, fill image opacity=1}]
                \includegraphics[width=\linewidth]{#3}
            \end{tcolorbox}
        \end{center}
        
        \begin{lowerpartauto}
            \RoleType{#2}
            \RoleObjective{#4}
            \RoleAbilities{#5}
            \VictoryCondition{#6}
            \StrategyTip{#7}
            \RoleWarnings{#8}
        \end{lowerpartauto}
    \end{neutralroleauto}
}

% Role type header
\newcommand{\RoleType}[1]{
    \textbf{\large #1 Role}\\[0.5ex]
}

% Role objective
\newcommand{\RoleObjective}[1]{
    \textbf{Objective: }#1\\[0.25ex]
}

% Role abilities
\newcommand{\RoleAbilities}[1]{
    \textbf{Special Abilities: }#1\\[0.25ex]
}

% Victory condition
\newcommand{\VictoryCondition}[1]{
    \textbf{Victory Condition: }#1\\[0.25ex]
}

% Strategy tip
\newcommand{\StrategyTip}[1]{
    \textbf{Strategy Tip: }\textit{#1}\\[0.25ex]
}

% Role warnings
\newcommand{\RoleWarnings}[1]{
    \ifx\relax#1\relax
    \else
        \textbf{Warning: }#1\\[0.25ex]
    \fi
}

% Conditional text helpers
\newcommand{\RoleTag}[2]{
    \textbf{#1: }#2\\[0.25ex]
}

\newcommand{\FlavorQuote}[1]{
    \begin{center}
    \textit{"#1"}
    \end{center}
    \\[0.25ex]
}

% Alignment-based coloring
\newcommand{\AlignmentColor}[2]{
    \ifnum#1=1
        \textcolor{blue}{\textbf{#2}}
    \fi
    \ifnum#1=2
        \textcolor{red}{\textbf{#2}}
    \fi
    \ifnum#1=3
        \textcolor{orange}{\textbf{#2}}
    \fi
    \ifnum#1=0
        \textbf{#2}
    \fi
} 